\section{統計の限界とバイアス : 手法の外側を知る}

\subsection{正しく計算しても結論が間違うとき}

前章の終わりに, 手法の外側を知ることが手法を正しく使うための最後の一歩になる, と書いた.
その外側で何が起こるのかは, 1つの実例を見るのが早い.

1973年, カリフォルニア大学バークレー校の大学院は, 性別を理由に女子受験生を差別しているとして訴えられた.
証拠とされたのは合格率の差だ. 全志願者を集計すると, 男子の合格率は約44\%, 女子は約35\%. 差は9ポイントある.
この数字を見て, 大学院に問題があると判断した人は当時多かった.

ところが大学院側が学部別に集計し直してみると, 数字の様子が変わる.
6つの主要学部のうち4つで, 女子の合格率は男子と同じか, むしろ高かった.
残り2つも, 全体集計に出ていた9ポイント差ほど大きな男子有利ではない.
学部単位で見るかぎり, 「女子の方が合格しにくい」と言える証拠はどこにも見当たらない.
それなのに, 大学全体の合計だけが男子有利に出る.

なぜそんなことが起こるのか.
学部ごとの合格率が, そもそも大きく違っていた.
人文系の合格率は20\%から30\%, 理工系は60\%から70\%という具合に, 学部間で2倍3倍の差がある.
そして女子は前者の「狭き門」に多く志望し, 男子は後者の「広き門」に多く志望していた.
各学部の中では公平でも, 女子が狭き門の学部に集中していたために, 全体平均では女子の合格率が下に引かれる.

集計の単位を変えただけで結論が逆転する. このような現象を\textbf{シンプソンのパラドックス}と呼ぶ.
平均は2章で, 割合は5章で扱った. どの計算も正しい.
それでも, 全体集計と学部別の集計で正反対のことが言える.

ただし, 集計の単位そのものが結論を狂わせたのではない.
学部によって合格率が大きく違い, 学部選びに男女差があった. この2つが重なって, 「学部」という第三の変数が性別と合格率の両方に絡んでいた.
全体集計はその第三の変数を隠し, 学部別の集計は表に出す. 結論を分けたのは, 集計の単位の裏で隠れたり現れたりしていた第三の変数のほうだ.
群によって混ざり方の違う第三の変数がある場面なら, 同じ逆転は平均の比較でも割合の比較でも起こりうる.

2章から9章までの手法は, 手順どおりに計算すれば誰がやっても同じ数字を返す. それでも結論が間違うのなら, 歪みは計算の外から入っている.
1章で述べたとおり, 標本には母集団のすべての情報が入らない. 特定の対象への偏りが生じるかどうかは, データの集め方による. だからこそ, 計算結果だけでなく, どのように集めたデータなのかも点検しなければならない. 本章では, 集め方から解釈までに入り込む歪みを, データの流れに沿って上流から順にたどる.
標本をどう集めたか. 集まったデータの中で何が観測できていたか. 出てきた数字をどう解釈したか.
シンプソンのパラドックスは, このうち最後の解釈の段階で起きた歪みにあたる. 集計の単位という解釈上の選択が, 結論を逆転させていた.
歪みの3類型を整理した後, 手法そのものが暗黙に置いている前提, そして数字の読み方の誤りへと進み, 最後に本書全体を振り返って閉じる.


\subsection{バイアスの3類型 : 上流・中流・下流の歪み}

3類型のうち最も身近なのは上流, 標本の集め方に入る歪みだ.

\subsubsection*{上流 : サンプリングバイアス}

5章の終わりに, \textbf{精密に間違える}という言葉が出てきた.
標本の集め方が偏っていると, $n$を増やすほど信頼区間は狭くなるのに, 区間の中心は的外れのままになる, という警告だった.
あの警告に, ここで名前をつける.

たとえば, 関心のある母集団は「日本に住む有権者」なのに, 集まったのが「平日の昼に渋谷駅にいた人」だったとしよう.
件数を10万人に増やしたところで, 平日の昼に渋谷駅にいない人の意見は最後までゼロのままだ.
標本平均が近づいていく先は, 知りたい「日本に住む有権者」の平均ではなく, 「平日の昼に渋谷駅にいた人たち」の平均でしかない.

これが\textbf{サンプリングバイアス}である.
標本が母集団を代表していないとき, 信頼区間は$n$を増やすほど狭く, しかも誤った位置に決まっていく.
件数を増やして解決できる問題ではない.

5章で紹介したリテラリー・ダイジェスト誌の大はずれも, 構図は渋谷駅と同じだった.
240万人という件数はギャラップの5万人を数十倍上回っていたのに, 電話帳と自動車登録簿から作った名簿が, 大恐慌下で電話も車も持たない貧困層をまるごと外していた.
ギャラップが勝ったのは件数ではなく, 標本の代表性で勝っていたからだ.

サンプリングバイアスの厄介さは, データの中身をいくら眺めても見抜けないところにある.
点検の対象は, 中身ではなく抽出の手順そのものだ.
確認するのは2点. 誰がリストに載っていたかという抽出枠と, 声を拾えた人と拾えなかった人を分ける回答率である.
ダイジェスト誌の調査なら, 「電話帳と自動車登録簿から抽出した」という1行に, 1936年の貧困層が入っていない事実がそのまま書かれている.
分析を始める前に, 抽出枠と回答率をまず確かめる. サンプリングバイアスの大半は, この入口の確認で見つかる.

\subsubsection*{中流 : 生存者バイアス}

抽出の枠組みの点検では捕まらない歪みもある.
「そもそも観測できなかったもの」の抜け落ちだ.

第二次大戦中, 米軍は帰還した爆撃機の被弾箇所を調べて, どこを補強すべきかを検討していた.
集計してみると, 翼の先端と機体後部に弾痕が集中していた.
最初の判断はこうだ. 「弾痕の多い場所を補強しよう」.

統計学者アブラハム・ワルドはこの結論に待ったをかけた.
補強すべきは弾痕の少ない場所のほうだ, と.
帰還した機は, 翼や機体後部に被弾しても飛び続けられた機だ.
エンジンや操縦席を被弾した機は, そもそも基地に戻ってこない.
分析の対象に上がっていたのは「生き残った機」だけで, 「生き残れなかった理由」がデータから抜けていた.

これが\textbf{生存者バイアス}である.
「観測できたものだけ」を見て結論を出すと, 観測できなかったものが体系的に抜け落ちる.
事業継続中の企業だけを集めて成功要因を語る研究, 在籍中の社員だけを集めた満足度調査, 完走したマラソン選手だけを集めた走法分析.
どれも, 退場した側に重要な情報が眠っている可能性がある.

点検の手順は, 手元のデータを見るより先に, 「ここに本当はどんな対象が含まれるべきか」を書き出すことだ.
書き出した母集団と実際に手元にあるデータの差が, 生存者バイアスの入る隙間にあたる.
その差を説明できないなら, 結論はその差の分だけ歪んでいる.

\subsubsection*{下流 : 確証バイアス}

3つめの歪みは, データではなく, 数字を読む人間の側に入る.
人間には, 自分の仮説に合うデータを重く受け取り, 合わないデータを軽く流す癖がある.
この癖を\textbf{確証バイアス}と呼ぶ.
心の持ちようの問題に聞こえるかもしれないが, 確証バイアスは統計の手続きの中にも入り込む.
選び方しだいで, 1つ1つの計算は正しいまま, 全体として嘘をつけてしまうからだ.
6章の検定を使って, その仕組みを具体的に見る.

公平なコインを20枚用意して, それぞれ100回ずつ投げ, 各コインで「表と裏の出方が偏っているか」を有意水準5\%で検定する場面を考えよう.
有意水準5\%とは, 6章で導入したとおり「本当は公平なのに偶然偏って見える確率を5\%まで許容する」という設定だ.
コイン1枚あたり, 偽陽性（公平なのに有意と判定される）が出る確率は5\%, 出ない確率は95\%である.
20枚を独立に検定するから, どの1枚も偽陽性を出さない確率は$0.95^{20}$. その裏返しで, どこかで偽陽性が出る確率は

$$
1 - (1 - 0.05)^{20} \approx 0.64
$$

つまり約64\%だ. 20回試せば, 6割以上の確率で「有意な発見」が偶然で生まれる.

ここで実験者が「有意になった1枚」だけを論文にし, 残り19枚を引き出しにしまったらどうなるか.
表に出るのは「特殊なコインを発見した」という主張だけだ.
6章で習ったp値の計算は, どの1枚についても正しく行われている.
1枚ごとの検定だけ見れば嘘はない.
それなのに, 「20枚試して当たりだけ報告した」という1段上の選別のせいで, 全体としての結論は嘘になる.

これが\textbf{p-ハッキング}と呼ばれる行為だ.
自分の仮説に合う結果だけを拾うという確証バイアスの癖が, ここでは検定という正式な手続きの上で動いている.
都合のいい結果が出るまで分析の角度や対象を変え続けると, 拾う作業に自覚すら要らなくなる.

抑える手順は2つ用意されている.
1つは事前登録. 何を, どんな仮説で, どの分析手順で見るかを, データを見る前に書いて第三者に渡しておく.
やってから決めるのではなく, やる前に決めてあるので, 「都合のいい結果が出てから条件を変える」という余地がふさがる.
もう1つは多重比較の補正. 同じデータで20回検定するなら, 各検定の有意水準を5\%でなく0.25\%程度に厳しくする（ボンフェローニ補正）か, 別の方法で「全体の偽陽性率」を抑える.
補正の詳細は本書の範囲を超えるが, 枠組みが用意されているので, 必要になった場面で調べ直せば足りる.

3類型は別々の知識ではなく, 標本の入手から解釈までの1本の流れの上の3地点である.
点検の時期も流れに沿う. 分析を始める前に抽出枠と回答率を確かめ, 手元のデータを前にしたら観測できなかったものを書き出し, 結果を読むときに選別の有無を疑う.


\subsection{手法の前提が崩れるとき : 独立性・線形性・等分散性}

データの集め方に歪みがなくても, 手法そのものは暗黙の前提に乗っている.
前章の終わりに, 最小二乗法と$R^2$の前提を3つ挙げた. 関係がおおむね直線的であること, データ点が互いに独立であること, 極端な外れ値がないこと.
このうち外れ値の怖さは, 1点が傾きを歪めるアンスコムの3組目として前章で既に見た.
ここでは残る2つ, 独立性と線形性に, 幅つきの話で効いてくる等分散性を加えた3つを, 回帰に限らず2章から9章の推論全体に関わる前提として点検する.
崩れていたときに何ができるかも, それぞれに添える.

\subsubsection*{独立性 : 隣の点と関係していないか}

4章のCLTの前提として「独立同分布」を導入した.
ここでの「独立」は, データ点どうしが互いに無関係であることを指す.
ある月のデータが翌月のデータに影響しているなら, それは独立ではない.

前章では製作費（説明変数）と視聴回数（目的変数）の作品データに対して回帰直線を引いた.
回帰の計算式は, 各データ点を独立に扱い, 残差の二乗和を最小化している.
ところが, データ点が独立でなかったらどうなるか.
たとえば同じシリーズの続編作品が複数本含まれていれば, それらの視聴回数は互いに似た値になりやすく, 隣どうしの残差は同じ符号をとりやすい.
時系列データの場合も同じで, ある月の値が翌月の値に影響していれば隣の残差が連動する.

残差が独立に動くとき, それぞれが上下にランダムに散らばるから, 残差全体の散らばりが「データの本当の不確かさ」をきちんと写し取る.
ところが残差が隣どうしで似ていると, 「実は独立に取り直せばもっと違う値が出るはずなのに, 連動して動いているせいで似た値しか観測されていない」状態になる.
散らばりの計算は観測されたものをそのまま反映するから, 連動を割り引いてはくれない.
結果として, 計算上のSEは本当の不確かさより小さくなる.

SEが小さいと信頼区間も狭く出るし, 6章のp値も小さく出る.
本当は不確かな結論が, 数字の上では確実そうに見えてしまう.

時系列データだけが該当するわけではない.
クラスごとの生徒, 同じ会社の社員, 同じ地域の住民. 集団内の点は互いに似ている傾向がある.
集団効果を無視して全員を独立だと扱うと, 同じ問題が起こる.

独立性が崩れていそうなときの対処は, 「独立でない構造を計算式の側に取り込む」ことだ.
時系列データには時系列モデル（AR・MAなど）が, 集団内で似ているデータには階層モデル（マルチレベル分析）が用意されている.
本書の範囲は「独立ではなさそうだ」と気づくところまでで, その先は手法の名前を頼りに調べがつく.

\subsubsection*{線形性 : 直線で表せる関係なのか}

直線モデルは「関係がおおむね直線で表せる」という前提の下で意味を持つ.
本当の関係が放物線, 指数, 周期的なものだったら, 直線を引いて得られるのは「もっとも近い直線」だけだ.
傾きや切片の数値は出るが, それらは現実の関係を要約していない.

前章で残差プロットを導入したとき, アンスコムの2組目の話をした.
4組のデータはどれも$r \approx 0.82$, $R^2 \approx 0.67$という似た数字を返したが, 2組目の実体はきれいな放物線で, 直線を当てはめること自体が不適切だった.

線形性が崩れているかどうかは, 残差プロットに目を向ければ見える.
予測値の小さい所と大きい所の両方で残差が正に偏り, 中央で負に偏っていれば, 関係はU字型だ.
逆向きのパターンも同じ警告にあたる.

線形性の確認は数値だけでは届かない.
$R^2$が0.7を超えていても, 残差プロットが系統的なパターンを示していれば直線モデルは合っていない.
逆に$R^2$が0.4でも, 残差にパターンがなければ, それはそれで意味のある直線関係だ.
$R^2$は当てはまりの「総量」しか語らず, 「外し方の質」は残差プロットでしかわからない.

崩れていたときの対処は, 関係の形に合わせて式を変えることだ.
$x^2$の項を加えて2次式にする. $\log x$を取って線形に直す. 周期成分があれば三角関数を加える.
線形回帰の計算自体は, 説明変数を増やしても同じ手順で実行できるから, 「直線でないなら捨てる」のではなく「形を変えて当てはめる」と思っておくと使いやすい.

\subsubsection*{等分散性 : 場所によって散らばりが違わないか}

3つめの前提が, 残差の散らばりがデータの位置によって変わらないこと, つまり等分散性である.
7章のt検定でWelch法を採ったのは, 2群間で分散が違うことを許す扱いを選んだという意味で, 等分散性の仮定を緩めた選択だった.

回帰では事情が少し違う.
たとえば, 予測値が大きい領域ほど残差の散らばりが広がるデータでは, 残差プロットに左側が狭く右側が広い扇形が現れる.
この扇形は, 等分散性が崩れたときの典型的な形の一つである.
第9章の残差プロットでも, 残差の幅が場所によって一定とは言いにくかった.
ただし, 大きな残差は予測値の小さい側にも現れており, 右へ広がる典型的な扇形ではなかった.
等分散性の崩れを疑う形は, 扇形だけではない.

このとき, 直線そのものは間違っていない.
平均的な傾きは正しく推定されている.
ただし, 「製作費を5億円に決めたら視聴回数はおよそいくつか」と幅つきで予測したい場合, 場所ごとに本当の散らばりが違うのに一律のSEで予測区間を出すと, 散らばりが小さい場所では広すぎ, 大きい場所では狭すぎる.
平均の推定は使えても, 予測の幅の保証は当てにならない. 「全部ダメ」ではなく「部分的にダメ」だ.

対処も独立性と線形性の場合と同じように用意されている.
残差の散らばりに合わせてSEを計算し直す方法（ロバスト標準誤差）と, 残差の大きい点を軽く扱う方法（重み付き最小二乗法）で, どちらも名前から調べがつく.

3つの前提の点検は, ほとんど残差プロット1枚に集約される.
残差を予測値に対してプロットし, ランダムに散らばっていればよし, 系統的なパターンや場所による幅の違いが見えたら警告.
前章で導入したあのプロットは, 線形回帰を使ってよいかどうかの確認手順を兼ねていたわけだ.


\subsection{数字の読み方を誤る典型 : p値・相関・外挿}

集め方にも手法の前提にも問題がないのに, 結論だけが間違う場面が残っている.
出てきた数字の読み方を誤る場合だ.
誤読が集中する場所は3か所に絞れる. p値, 相関と因果の関係, データ範囲の外側である.

\subsubsection*{p値の3つの誤読}

6章でp値を「帰無仮説が正しいと仮定したとき, 観測値以上に極端な結果が出る確率」と定義した.
この定義はそれなりに込み入っているせいで, 短く要約しようとすると意味がずれる.
ずれ方には型がある.

最初の誤読は, 「p値は帰無仮説が正しい確率である」という読み方だ.
これは違う. p値は「帰無仮説が正しいと仮定したとき」の珍しさを測っているのであって, 帰無仮説そのものの真偽の確率を返してくれるわけではない.
p値が0.03だからといって, 帰無仮説が正しい確率が3\%という意味にはならない.

次の誤読は, 「p<0.05なら効果は大きい」という読み方だ.
p値は珍しさの指標であって, 効果の大きさは別の指標で測る.
両者が別物である理由は, 「件数が増えると有意になりやすい」という性質に表れている.
4章で見たとおりSEは$u/\sqrt{n}$で, $n$が大きくなると小さくなる.
6章のt検定では検定統計量$t = \text{(2群の平均差)}/\text{SE}$だから, SEが小さくなれば$t$は大きくなる.
$t$が大きくなれば, 同じ平均差でもp値は小さく出る.
だから2群の平均差が小さくても, $n$が大きければ「珍しい」と判定される.
逆に$n$が小さければ, それなりに大きな差でも有意にならないことがある.

最後の誤読は, 「p>0.05なら効果はない」という読み方だ.
これも違う. 「珍しいと言える証拠が足りない」と「効果がないことが証明された」は別の話だ.
件数が少ない実験で有意にならなかったからといって, 効果が存在しないわけではない.

p値を扱うときは, p値だけで報告を済ませない方がよい.
5章の信頼区間と, 7章で触れた効果量を一緒に出す.
信頼区間は「効果がだいたいどの範囲にありそうか」, 効果量は「差そのものの大きさ」, p値は「偶然で説明しにくいかどうかの目安」.
3つはそれぞれ別の問いに答えている.

具体的な書き方の目安はこうだ.

\begin{quote}
2群の平均差は1.2点（95\%信頼区間 0.4から2.0）, 効果量（Cohenの$d$）は0.3, $p = 0.01$.
\end{quote}

平均差1.2点の周りに信頼区間が0.4から2.0と置かれているおかげで「効果は小さめだが正の方向は確かそうだ」と読み取れ, $d = 0.3$で「効果としては小さい部類」と相対化できる. p値はその仕上げに過ぎない.
報告するときはこの3点を並べる. 報告を読むときは, 3点が揃っているかをまず確かめる.

\subsubsection*{相関は因果ではない, ではどう違うのか}

相関が因果を意味しないことは, 8章で\textbf{擬似相関}と\textbf{交絡因子}という名前とともに確かめた.
アイスクリームの売上と水難事故の相関を作っていたのは, 両方を同時に動かす気温という第三の変数だった.
8章の終わりには, この点は第10章で改めて扱う, と書いた. その約束をここで果たす.

本章の冒頭のバークレーの話は, 交絡の実例そのものだ.
性別と合格率のあいだに「学部」という第三の変数が絡み, 直接の差別がない学部別の数字から, 全体では9ポイントの差が生まれていた.
気温のような物理量にかぎらず, 学部の選ばれ方, 国の面積や人口のような構造的な量まで, 両方の変数に絡むものは何でも交絡因子になりうる.
ヨーロッパ各国でコウノトリの巣の数と人間の出生数に正の相関が出るという有名な話も, 国の規模が両方を押し上げているだけである.

交絡を疑うときの確認手順は, 「両方を同時に押し上げそうな第三の変数」を1つ思いつくことから始まる.
思いついたら, その変数の値ごとにデータを分けて計算し直す. この手続きを\textbf{層別分析}と呼ぶ.
バークレーの大学院側がやったのは, まさにこの手続きだった. 学部で層別して合格率を出し直したら, 全体集計の男女差はほぼ消えた.
気温で層別したアイスと水難事故の相関も, 同じように消えるはずだ.
層別して消えれば交絡. 消えなければ, 別の説明を探すことになる.

その「別の説明」の筆頭が\textbf{逆因果}である.
$x$が$y$を動かしているのではなく, $y$が$x$を動かしている場合だ.
8章のおすすめ表示と視聴回数にも, この疑いがあった. 表示されたから視聴されたのか, よく視聴されたから表示が増えたのか.
「医療支出が多い地域ほど健康指標が悪い」というデータでも同じことが起こる.
医療支出が健康を損ねているのではなく, 健康状態が悪い地域だから医療支出が増えている, という向きの因果かもしれない.
矢印の向きを取り違えると, 結論はそのまま逆転する.

逆因果を疑うときの目印は時間順序だ.
$x$の変化は$y$の変化より先に起きているか.
データが同じ時点のスナップショットなら, 因果の向きは原理的に決まらない.
別の時点で同じ現象を測れたときに初めて, 矢印の向きを絞り込む手がかりが得られる.

交絡も逆因果もないのに強い相関が出ることもある. \textbf{偶然の一致}である.
無関係な2つの量を大量に集めて相関を計算すれば, そのうちのいくつかは偶然強く相関する.
「米国のチーズ消費量とベッドシーツでの絞死者数」のような, ふざけた相関の実例集が世間に出回っているくらいだ.
たくさん試せば偽陽性は出る. 本章の前半で見たp-ハッキングと同じ構造が, 相関の探索でも起こる.

偶然か本物かを切り分ける最低限の手は, 別のデータで再現するかを確かめることだ.
別の標本でも同じ強さの相関が出るなら本物の見込みが高い. 1度きりなら, 偶然の可能性は消えない.

交絡・逆因果・偶然は, 1つずつでも併発でも起こる.
相関を見たら問うことは3つ. 両方を押し上げる第三の変数はないか. 因果の向きは逆ではないか. 似たような量を大量に試した中の1つではないか.
1つでも即答できないうちは, 「強い相関がある」というだけで結論を下せない.

\subsubsection*{データの外側に踏み込むとき}

3つめの誤読は, 前章の終わりで一度通った場所にある.
製作費を億円, 視聴回数を万回で表した回帰直線$y \approx 108.4 + 49.4x$に製作費30億円を入れると, 視聴回数は約1590万回と計算される.
しかしデータの範囲は0.25億円から11.46億円で, そこで観測された直線関係が30億円まで続いている保証はどこにもない.
一定の額を超えると, 制作やプロモーションの仕組み自体が変わる場合もある.
範囲の外に直線を延ばす外挿は, 「直線が続いている」という根拠のない仮定を黙って置く行為なのだった.
切片$a \approx 108.4$を「製作費ゼロでも視聴回数108万回」と読めないのも同じ理由で, 製作費ゼロの付近にデータはなく, 108.4は直線の位置決めの計算から出た数値に過ぎない.

前章では回帰の話として出てきたが, 誤読の型として見ると, 外挿は回帰の直線に限らない.
過去3年の売上の伸び率を来年にそのまま延ばす. 20代への調査結果を全年代に広げて語る.
どちらも, 観測した範囲で成り立っていた関係を, 観測のない場所へ黙って持ち出している.

確認の手順は2つ.
語ろうとしている点が観測範囲の中にあるかを確かめる. 範囲の外を語るなら, 外挿であることを明記する.
数字が出ることと, その数字を現実に当てはめてよいことは別である.
モデルが関係を要約しているのは, 観測が届いた範囲までだ.


\subsection{本書の終わりに}

本書を1章から振り返ると, ずいぶん遠くまで来た.
要約から始めて, データの形を見て, ばらつきから幅を計算し, 群の差を吟味し, 関係を測り, モデルを当てはめ, そして本章で手法の外側にある歪みの見取り図を描いた.
学び始めの頃は, 計算結果を正しく出すことが目標だったかもしれない.
振り返ってみると, 正しく計算するのは出発点で, 本書の後半はずっと, 出てきた数字を疑う手順の話をしていた.
出した数字の前提は何か. 集計の単位は妥当か. 第三の変数を見落としていないか. 観測範囲の外まで延ばして語っていないか.
数字を出すたびにこの問いへ戻ってもらえるなら, 本書の役目はおおむね果たされている.

本書は統計の骨格に絞ったので, 各章で名前だけ挙げた話題には, それぞれ教科書1冊分の世界が広がっている.
多重比較の制御, ブートストラップ, 重回帰, 時系列モデルや階層モデルといった話題だ.
どの章で「もっと知りたい」と感じたか. その章の続きから手を伸ばすのが自然な順序だ.

最後にひとつ書いておきたい.
本書の手順を, 自分の手元のデータで, 最初から最後まで一度回してみてほしい.
その途中で引っかかった「ここはなぜそうなのか」が, 統計学の次の章への入口になる.

ここでいったん筆を置く. 続きは, 読者自身のデータの上で書かれる.
